\begin{enumerate}
\item \textbf{Data selection} (selecting appropriate images showing the
  evolution of the CME) and identification of CME front.

  % FIGURE NEEDED: mosaic of 7 JHelioviewer frames (SDO/AIA + SOHO/LASCO
  % coronagraph images) numbered 1--7, each with the CME front marked by a
  % line from the solar limb; 2021_SOLAR.pdf, p.~1.

  A total of 7 time intervals are selected tracing the evolution of the CME:
  \begin{enumerate}
  \item 2021/05/28 22:19:00 UT
  \item 2021/05/28 22:54:11 UT
  \item 2021/05/28 23:30:08 UT
  \item 2021/05/29 23:42:08 UT
  % sic: the official solution lists 2021/05/29 23:42:08 UT here, while its
  % summary table gives 2021/05/28 23:42:08 UT for interval 4
  \item 2021/05/29 01:12:16 UT
  \item 2021/05/29 03:36:08 UT
  \item 2021/05/29 05:00:08 UT
  % sic: the summary table gives 2021/05/29 05:36:08 UT for interval 7, but
  % 05:00:08 UT is consistent with the quoted 401 minutes from onset
  \end{enumerate}

  Summary of the information extracted from the images for all different
  selected intervals:

  {\footnotesize\setlength{\tabcolsep}{2.5pt}
  \begin{center}
  \begin{tabular}{c l c c c c c c}
  \toprule
  Interval & Date and Time & \begin{tabular}{@{}c@{}}Time\\ interval\\ (min)\end{tabular}
    & \begin{tabular}{@{}c@{}}Time from\\ CME onset\\ (min)\end{tabular}
    & \multicolumn{2}{c}{\begin{tabular}{@{}c@{}}Coordinates of\\ the CME front\end{tabular}}
    & \begin{tabular}{@{}c@{}}Distance from\\ disc center\\ (arcsec)\end{tabular}
    & \begin{tabular}{@{}c@{}}Distance from\\ disc center\\ (km)\end{tabular}\\
  & & & & $x''$ & $y''$ & &\\
  \midrule
  1 & 2021/05/28 22:19:00 UT & 0   & 0   & 950   & 144     & 961   & 696617\\
  2 & 2021/05/28 22:54:11 UT & 35  & 35  & 2682  & 174     & 2688  & 1948538\\
  3 & 2021/05/28 23:30:08 UT & 36  & 71  & 4579  & 97      & 4580  & 3320520\\
  4 & 2021/05/28 23:42:08 UT & 47  & 83  & 7441  & 116     & 7442  & 5395380\\
  5 & 2021/05/29 01:12:16 UT & 126 & 173 & 14000 & $-891$  & 14028 & 10170535\\
  6 & 2021/05/29 03:36:08 UT & 191 & 317 & 21048 & $-4570$ & 21538 & 15615349\\
  7 & 2021/05/29 05:36:08 UT & 210 & 401 & 27922 & $-7126$ & 28817 & 20892306\\
  \bottomrule
  \end{tabular}
  \end{center}
  % sic: the "Time interval" column of the official table is not consistent
  % with the differences of the successive timestamps (rows 4--7)

  \begin{center}
  \begin{tabular}{c c c c c c c}
  \toprule
  Interval & \begin{tabular}{@{}c@{}}Distance\\ from onset\\ (km)\end{tabular}
    & \begin{tabular}{@{}c@{}}Cumulative\\ velocity, plane\\ of the sky (km/s)\end{tabular}
    & \begin{tabular}{@{}c@{}}Deprojection of\\ velocity towards\\ ParkerSP (km/s)\end{tabular}
    & \begin{tabular}{@{}c@{}}Distance per\\ time interval\\ (km)\end{tabular}
    & \begin{tabular}{@{}c@{}}Velocity on\\ plane of the\\ sky (km/s)\end{tabular}
    & \begin{tabular}{@{}c@{}}Deprojection of\\ velocity towards\\ ParkerSP (km/s)\end{tabular}\\
  \midrule
  1 & 0        & 0   & 0    & 0       & 0   & 0\\
  2 & 1252198  & 596 & 728  & 1252198 & 596 & 728\\
  3 & 2624180  & 616 & 752  & 1371982 & 635 & 775\\
  4 & 4699040  & 944 & 1152 & 2074861 & 736 & 898\\
  5 & 9474195  & 913 & 1114 & 4775155 & 632 & 771\\
  6 & 14919009 & 784 & 958  & 5444814 & 475 & 580\\
  7 & 20195966 & 839 & 1025 & 5276957 & 419 & 511\\
  \bottomrule
  \end{tabular}
  \end{center}
  }
\item \textbf{Plotting the kinetic evolution of the CME.}

  Plots: distance vs time and velocity vs time.

  % FIGURE NEEDED: three plots -- Distance vs Time, Cumulative Velocity vs
  % Time, and Velocity (per interval) vs Time, built from the table above;
  % 2021_SOLAR.pdf, p.~2.
\item \textbf{Calculating the velocity of the CME front reaching ParkerSP.}

  The velocity corresponding to the last interval (up to reaching 30 solar
  radii) is 511\,km/s which, according to the considerations of the problem,
  would be the same velocity of the CME front once it reaches the ParkerSP.

  \textbf{Calculating the time spent by the CME front to reach the ParkerSP
  from its onset.}

  The Sun--ParkerSP distance estimated from Figure 1 is approx.\
  $0.66\,\mathrm{AU} = 9.9\times10^{7}\,\mathrm{km} = 141$ solar radii (given
  the information that 10 solar radii is approx.\ $7\times10^{6}$\,km).

  From the previous analysis, it is known that the CME front reaches a
  distance of 30 solar radii after 401 minutes (6.7 hours) from its onset
  (see Table). The distance that the CME travels at the constant speed of
  511\,km/s is expressed as $141-30=111$ solar radii
  $= 7.8\times 10^{7}$\,km, and thus the time spent in that interval is
  42.4 hours.

  Finally, the total time it takes for the CME front to reach the ParkerSP
  from its onset is computed as: $6.7 + 42.4 = 49.1$ hours.
\item \textbf{True and false statements}
  \begin{description}
  \item[A.] A larger number of selected images would give more precise
    information on the evolution of the CME and the calculated physical
    parameters. TRUE
  \item[B.] The analysis and measurements of the CME evolution should
    consider the differential rotation of the Sun, and therefore the
    calculated velocities will be affected. FALSE
  \item[C.] Any software (numerical) misalignment among the images when
    creating the mosaic will have direct effects on the precision of your
    calculations. TRUE
  \item[D.] The different assumptions made in order to construct the model
    displaying in a heliospheric density map in Figure 1, may affect the
    estimation of the Sun--ParkerSP distance. FALSE
  \item[E.] The interaction of the CME-front with the remnant dust left by
    the 2019 Borisov comet broadens and difuminate the images. This
    intensifies the lack of contrast in the images, increasing the
    uncertainty in determining the CME-front and its propagation. FALSE
    % sic: "difuminate" as in the original
  \end{description}
\item \textbf{Kinetic energy calculations}

  Proton:
  \[
  0.5 \times (1.67262\times10^{-27}~\mathrm{kg}) \times (511~\mathrm{m/s})^2
  = 1.4\times10^{-3}~\mathrm{eV}
  \]
  Alpha particle:
  \[
  0.5 \times (6.64\times10^{-27}~\mathrm{kg}) \times (511~\mathrm{m/s})^2
  = 5.4\times10^{-3}~\mathrm{eV}
  \]
  % sic: the official solution uses 511 m/s (instead of the CME speed of
  % 511 km/s); the quoted eV values follow from the m/s figure.
\end{enumerate}
